% GENERATED FILE. Edit canonical lessons, then run npm run generate:books.
% canonical-source-hash: fnv1a64:90abd5420289c825
% canonical-lessons: ZH-C131-chuntian, ZH-W131-chun-spring, ZH-C131-xiatian, ZH-C131-qiutian, ZH-C131-dongtian, ZH-C131-jijie

\chapter{Four Seasons}
\label{ch:zh-four-seasons}
\hlchaptermodality{131}

\begin{chapteropening}
\textbf{By the end of this chapter:} \emph{I can name the four seasons and ask which one someone likes.}

Complete the last lesson of chapter 131: I can name the four seasons and ask which one someone likes.
\end{chapteropening}

\section[\texorpdfstring{\zh{春天} (chūntiān)}{春天 (chūntiān)}]{\zh{春天} (chūntiān) — spring}

\label{lesson:ZH-C131-chuntian}



\begin{quote}
Before the new one: say the Mandarin for sore, then the Mandarin for to see a doctor.
\end{quote}

\subsection*{You'll want to know: \zh{春天}}
\textbf{\zh{春天}} — \emph{chūntiān} — \textquotedblleft{}spring\textquotedblright{}.

\begin{quote}
\textbf{\zh{我喜欢春天}} \emph{wǒ xǐhuan chūntiān} — I love spring.
\end{quote}

The new character, in pieces: \zh{春} — the three bars of \zh{三}, a \zh{人} whose two strokes spread across them, and \zh{日} underneath. The next lesson puts it in your hand.

\begin{cousinweb}
\textbf{\zh{春天}} (\emph{chūntiān}) is built from \zh{春} 'spring' and \zh{天} 'sky, day', both inherited from Middle Chinese, \emph{tsyhwin} and \emph{then} (Guangyun, in Baxter's transcription); the dictionary glosses the first 'the head of the four seasons'.
\end{cousinweb}

\begin{grammarlens}[title={What you've built}]
One more word for times of day.
\end{grammarlens}

\subsection*{Your turn}
\begin{itemize}
\raggedright
  \item \emph{Say it:} \emph{chūntiān}
  \item \emph{Say it:} \emph{chūntiān}, once more
  \item \emph{Say it:} \emph{wǒ xǐhuan chūntiān} — I love spring.
  \item \emph{Recall:} say the Mandarin for sore, then the Mandarin for to see a doctor, then say \emph{chūntiān} again
\end{itemize}

\begin{tcolorbox}[breakable,colback=teal!4,colframe=teal!35!black,title={Before you move on}]
What does \zh{春天} mean? (\textquotedblleft{}Spring\textquotedblright{}.) Say it once more.
\end{tcolorbox}
\section[\texorpdfstring{\zh{春} (chūn)}{春 (chūn)}]{\zh{春} — \zh{三} and \zh{人} over \zh{日}}

\label{lesson:ZH-W131-chun-spring}



\begin{quote}
The word first, then the character inside it.

\begin{itemize}
\raggedright
  \item \emph{Recall:} say \emph{chūntiān} — \textbf{R1}, one lesson back, before it decays
\end{itemize}

\begin{quote}
\zh{三} \zh{人} \zh{日}
\end{quote}

\emph{Read it:} each one aloud, with nothing said first — \emph{sān}, \emph{rén}, \emph{rì}
\end{quote}

\subsection*{Script — \zh{三} and \zh{人} over \zh{日}}
\hlblockfigure{\detokenize{figures/ZH-W131-chun-spring-filmstrip.pdf}}{How \zh{春} is written, stroke by stroke}

\begin{quote}
\zh{春}
\end{quote}

\textbf{Three pieces you know, top to bottom, the first two woven together.}

\textbf{The three bars of \zh{三}}: two shorter bars, then a long third one.

\textbf{\zh{人}, two strokes}, spread wide across the bars: the left-falling stroke from above the top bar down to the lower left, then the right-falling stroke from the middle out to the lower right.

\textbf{\zh{日}, four strokes}, small, under the spread of \zh{人}.

3 + 2 + 4 = nine.

\subsection*{Your turn}
\begin{itemize}
\raggedright
  \item \emph{Write it:} \zh{春}, all nine, with the model in front of you
  \item \emph{Cover:} write it again from memory, and count the pieces, not the strokes
\end{itemize}

\begin{tcolorbox}[breakable,colback=teal!4,colframe=teal!35!black,title={Before you move on}]
How many strokes in \zh{春}? (\textbf{Nine}.) Which bar is the longest? (\textbf{The third}.) What sits at the bottom? (\textbf{\zh{日}}.)

Source: \href{https://raw.githubusercontent.com/chanind/hanzi-writer-data/68d10a4b21150cae5e1ebbd223eed289cf32d90c/data/\%E6\%98\%A5.json}{Hanzi Writer Data \zh{春}, pinned stroke paths and medians}.
\end{tcolorbox}
\section[xiàtiān]{xiàtiān — summer}

\label{lesson:ZH-C131-xiatian}



\begin{quote}
Before the new one: say the Mandarin for to see a doctor, then the Mandarin for spring.
\end{quote}

\subsection*{You'll want to know: xiàtiān}
\emph{xiàtiān} — \textquotedblleft{}summer\textquotedblright{}.

\begin{quote}
\emph{Běijīng de xiàtiān hěn rè} — Summers in Beijing are hot.
\end{quote}

Its characters come later; for now it is a word for the ear.

\begin{cousinweb}
\emph{xiàtiān} is built from \emph{xià} 'summer' and \zh{天} 'sky', both inherited from Middle Chinese, \emph{haeH} and \emph{then} (Guangyun, in Baxter's transcription); the dictionary lists the first under the words 'spring and summer'.
\end{cousinweb}

\begin{grammarlens}[title={What you've built}]
One more word for times of day.
\end{grammarlens}

\subsection*{Your turn}
\begin{itemize}
\raggedright
  \item \emph{Say it:} \emph{xiàtiān}
  \item \emph{Say it:} \emph{xiàtiān}, once more
  \item \emph{Say it:} \emph{Běijīng de xiàtiān hěn rè} — Summers in Beijing are hot.
  \item \emph{Recall:} say the Mandarin for to see a doctor, then the Mandarin for spring, then say \emph{xiàtiān} again
\end{itemize}

\begin{tcolorbox}[breakable,colback=teal!4,colframe=teal!35!black,title={Before you move on}]
What does xiàtiān mean? (\textquotedblleft{}Summer\textquotedblright{}.) Say it once more.
\end{tcolorbox}
\section[qiūtiān]{qiūtiān — autumn}

\label{lesson:ZH-C131-qiutian}



\begin{quote}
Before the new one: say the Mandarin for spring, then the Mandarin for summer.
\end{quote}

\subsection*{You'll want to know: qiūtiān}
\emph{qiūtiān} — \textquotedblleft{}autumn\textquotedblright{}.

\begin{quote}
\emph{qiūtiān bù lěng yě bú rè} — Autumn is neither cold nor hot.
\end{quote}

Its characters come later; for now it is a word for the ear.

\begin{cousinweb}
\emph{qiūtiān} is built from \emph{qiū} 'autumn' and \zh{天} 'sky', both inherited from Middle Chinese, \emph{tshjuw} and \emph{then} (Guangyun, in Baxter's transcription); the dictionary quotes an older book for the first: the time when grain is ripe.
\end{cousinweb}

\begin{grammarlens}[title={What you've built}]
One more word for times of day.
\end{grammarlens}

\subsection*{Your turn}
\begin{itemize}
\raggedright
  \item \emph{Say it:} \emph{qiūtiān}
  \item \emph{Say it:} \emph{qiūtiān}, once more
  \item \emph{Say it:} \emph{qiūtiān bù lěng yě bú rè} — Autumn is neither cold nor hot.
  \item \emph{Recall:} say the Mandarin for spring, then the Mandarin for summer, then say \emph{qiūtiān} again
\end{itemize}

\begin{tcolorbox}[breakable,colback=teal!4,colframe=teal!35!black,title={Before you move on}]
What does qiūtiān mean? (\textquotedblleft{}Autumn\textquotedblright{}.) Say it once more.
\end{tcolorbox}
\section[dōngtiān]{dōngtiān — winter}

\label{lesson:ZH-C131-dongtian}



\begin{quote}
Before the new one: say the Mandarin for summer, then the Mandarin for autumn.
\end{quote}

\subsection*{You'll want to know: dōngtiān}
\emph{dōngtiān} — \textquotedblleft{}winter\textquotedblright{}.

\begin{quote}
\emph{dōngtiān wǒ bù xiǎng chūmén} — In winter I don't feel like going out.
\end{quote}

Its characters come later; for now it is a word for the ear.

\begin{cousinweb}
\emph{dōngtiān} is built from \emph{dōng} 'winter' and \zh{天} 'sky', both inherited from Middle Chinese, \emph{towng} and \emph{then} (Guangyun, in Baxter's transcription); the dictionary glosses the first 'the last of the four seasons'.
\end{cousinweb}

\begin{grammarlens}[title={What you've built}]
One more word for times of day.
\end{grammarlens}

\subsection*{Your turn}
\begin{itemize}
\raggedright
  \item \emph{Say it:} \emph{dōngtiān}
  \item \emph{Say it:} \emph{dōngtiān}, once more
  \item \emph{Say it:} \emph{dōngtiān wǒ bù xiǎng chūmén} — In winter I don't feel like going out.
  \item \emph{Recall:} say the Mandarin for summer, then the Mandarin for autumn, then say \emph{dōngtiān} again
\end{itemize}

\begin{tcolorbox}[breakable,colback=teal!4,colframe=teal!35!black,title={Before you move on}]
What does dōngtiān mean? (\textquotedblleft{}Winter\textquotedblright{}.) Say it once more.
\end{tcolorbox}
\section[jìjié]{jìjié — season}

\label{lesson:ZH-C131-jijie}



\begin{quote}
Before the new one: say the Mandarin for autumn, then the Mandarin for winter.
\end{quote}

\subsection*{You'll want to know: jìjié}
\emph{jìjié} — \textquotedblleft{}season\textquotedblright{}.

\begin{quote}
\emph{nǐ xǐhuan nǎge jìjié} — Which season do you like?
\end{quote}

Its characters come later; for now it is a word for the ear.

\begin{cousinweb}
\emph{jìjié} is built from \emph{jì} and \emph{jié}, both inherited from Middle Chinese, \emph{kjwijH} and \emph{tset} (Guangyun, in Baxter's transcription). The dictionary glosses the first 'the youngest of brothers' and the second, after an older book, 'the joint of a bamboo'; how the pair came to mean a season is not settled.
\end{cousinweb}

\begin{grammarlens}[title={What you've built}]
One more word for times of day.
\end{grammarlens}

\subsection*{Your turn}
\begin{itemize}
\raggedright
  \item \emph{Say it:} \emph{jìjié}
  \item \emph{Say it:} \emph{jìjié}, once more
  \item \emph{Say it:} \emph{nǐ xǐhuan nǎge jìjié} — Which season do you like?
  \item \emph{Recall:} say the Mandarin for autumn, then the Mandarin for winter, then say \emph{jìjié} again
\end{itemize}

\begin{tcolorbox}[breakable,colback=teal!4,colframe=teal!35!black,title={Before you move on}]
What does jìjié mean? (\textquotedblleft{}Season\textquotedblright{}.) Say it once more.
\end{tcolorbox}
